% =====================================================================
% paper/thesis_ch5/sections/intro.tex
% §5.1 章级独立引言（dissertation 体裁层）
% 起草依据：draft_5.1_outline.md（rev. 2026-06-06）+ spec rev.5 §0.5 体裁规范。
% 约束：全章自足、零跨章引述；抽象到"环境流场 + 单点局部可观测"，去环境/方法术语；
%       贡献定位为科学洞见而非新算法；online/offline 不对称权重；宽问题、窄证据；
%       负面结果作 finding；引言不携带实验数字（点到为止）。
% =====================================================================

\section{引言}\label{sec:ch5_intro}

自主航行器要在真实环境中可靠导航，必须应对一个自身无法完整观测的环境流场：流场在空间上非均匀、随时间演化，持续作用于整个载体而改变其运动，而航行器在部署时通常只能在自身所在的一点感知这一流场。这使得导航在本质上成为一个在不完全可观测的外部力场中持续决策的问题：驱动动力学的是作用于整体载体的流场，可获得的却只是局部、瞬时的一瞥。更关键的是，这种信息受限是部署现实强加的，而非建模选择：在真实海洋作业中，外部观测系统或高保真仿真所能提供的流场全貌均不可得，可依赖的只有载体自身携带的有限传感。学习式方法近年已被用于在复杂、非定常的流场中导航并展现出超越固定策略的潜力~\citep{gunnarson2021ncomm,verma2018pnas}，但它们在何种感知与数据条件下真正可行，仍缺乏系统刻画。

本章因此提出一个直接的问题：在只能依赖局部感知的设定下，强化学习能把基于局部信息的导航推进到什么程度？这一问题有两个面向：当航行器能与环境持续在线交互时，它能否仅凭局部感知学会导航、又是什么在限制它；以及——本章的核心——当只能从既有控制器留下的日志离线学习~\citep{levine2020offline}、且部署时不允许任何特权信息（仅训练阶段可得、部署时不可得的、对真实流场的额外了解）时，策略能逼近怎样的上界。后一设定更贴近部署现实：受限平台往往无法安全或廉价地大量在线试错，只能复用既有数据。本章据此先以在线情形确立任务的可学性与瓶颈所在，再以此框定并支撑真正的核心贡献——离线情形下的可行性。

本章给出的回答有三层，其中前两层与直觉相悖。其一，离线情形下，仅凭简单可部署控制器留下的日志，学到的策略即可追平一个被额外告知真实流场的特权基线，而部署时完全无需那条特权信息。其二，闭合不同感知配置之间的性能差距，靠的并非``看得更广''——增设空间上的传感探头——而是``如何随时间用好手头这一点信号''。其三，上述结论只在一定环境难度内成立：越过某一阈值后，会出现一个所有方法、包括被直接给定特权信息者都无法突破的根本感知极限，而且换用表达力更强的学习算法也并不普遍更优。

需要强调，这些结论是在一个受控的流场导航实验台上确立的；其精确边界——环境难度、传感配置与任务设定——将在后文逐一界定，本章并不将其包装为对一切流场或平台的普适结论。

与许多强化学习工作不同，本章的贡献不是一个新的学习算法，而是关于``什么决定成败''的实证洞见。所采用的算法与训练协议——在线 off-policy actor-critic、行为约束的离线方法、生成式策略先验，以及非对称 actor-critic 下的特权信息利用——均为已有成果~\citep{fujimoto2021td3bc,tarasov2023rebrac,kostrikov2022iql,pinto2017asymmetric}；本章的新意在于把它们置于同一受控的部署设定下，以统一的评估协议与多种子统计检验，得到一组可相互印证的发现：
\begin{itemize}
  \item 确立局部信息导航的可学性，并将其瓶颈定位于对部署信号的\emph{时序访问}，而非空间传感的丰富度；
  \item 证明仅凭可部署控制器的日志即可逼近特权基线、且部署无需任何特权信息，给出一条不依赖任何仿真专属信号即可逼近在线教师的配方；
  \item 刻画该配方的\emph{失效边界}：越过环境难度阈值后出现根本感知极限，连特权示范与特权 critic 也无法挽救；
  \item 表明算法优劣随数据与环境条件而变：真正的决定因素是模仿目标的质量与数据特性的匹配，而非策略先验的表达力。
\end{itemize}
贯穿这四点的，是两条正向杠杆——策略对部署信号的时序访问，以及模仿目标质量与数据条件的匹配。

本章其余部分按两幕组织，对应上述问题的两个面向。正文首先梳理相关工作并界定本章定位，再给出三条实验线共享的统一问题设定。第一幕处理在线情形，确立任务的可学性并将瓶颈定位于时序信息访问；第二幕处理离线情形，也是本章重心：先以一个可部署的离线基线拆解问题设定，继而由主线方法给出核心可行性结果，随后刻画其失效边界，最后通过算法对比揭示决定成败的真正因素。全章最后统一讨论贯穿各线的机制并对全章作结。
